\documentclass[12pt]{article}
\usepackage[utf8]{inputenc}
\usepackage[T1]{fontenc}
\usepackage{hyperref}
\title{Punthamara Language Resources for K-12 Curriculum Development}
\author{Compiled by Claude Sonnet 5.5 (claude-sonnet-5-5) for Andr\'as Kornai}
\date{}
\begin{document}
\maketitle

\section{Introduction}
Punthamara (ISO 639-3 xpt) is an extinct Aboriginal language of western Queensland, Australia. It is classified as Pama-Nyungan. It is listed as extinct since 1999, and the sources describe it as poorly documented. The main linguistic source is Nils Holmer's work (a 1979 manuscript and \emph{Notes on some Queensland languages}, 1988). Alternative spellings include Bunthomarra, Buntamara and Boonthamurra. One recorded word is \emph{ngaka} `water'. Any revitalization work should be led by the Punthamara descendants and by the relevant Aboriginal community bodies. No such body was identified in the search.

\section{Primary Educational Resources}
No primers, textbooks or school programmes were found.

\section{Linguistic and Reference Materials}
\begin{itemize}
\item Wikipedia: \url{https://en.wikipedia.org/wiki/Punthamara}
\item AIATSIS collection record: \url{https://collection.aiatsis.gov.au/node/1508862}
\item ASJP word list: \url{https://asjp.clld.org/languages/PUNTHAMARA}
\item PHOIBLE: \url{https://phoible.org/languages/punt1240}
\item Glottolog reference to R. H. Mathews (1905): \url{https://glottolog.org/resource/reference/id/584391}
\item Biodiversity Heritage Library part: \url{https://biodiversitylibrary.org/part/212925}
\end{itemize}

\section{Religious and Community Resources}
\begin{itemize}
\item South Australian Museum archives entry for the Punthamara language group: \url{https://www.samuseum.sa.gov.au/collection/archives/language_groups/punthamara}
\item State Library of Queensland blog, 2019 International Year of Indigenous Languages: \url{https://www.slq.qld.gov.au/blog/2019-international-year-indigenous-languages-word-week-week-six}
\end{itemize}
I could not verify what the last item says about Punthamara, because the search snippet did not show it.

\section{Digital and Audio Resources}
No audio or software resources were found.

\section{Curriculum Development Context}
Only wordlists and archival notes survive, so a curriculum could be only a small heritage unit. The first step is to contact AIATSIS, the State Library of Queensland and local Aboriginal communities about the manuscripts. Community consent and leadership are essential.
\end{document}